\documentclass[english, parskip=full]{scrartcl}
\usepackage[utf8]{inputenc}  % Kodierung der Datei
\usepackage[english]{babel}  % Sprache auf Deutsch 
\usepackage{color}
\usepackage{graphicx, gensymb}
\usepackage{amsfonts, cancel, amssymb}
\usepackage{amsmath, amsthm, array, esint}
\usepackage{booktabs, multirow}  % schönere Tabellen
\usepackage{caption}  % Anpassung der Captions
\usepackage{scrlayer-scrpage}    % Manipulation des Seitenstils
\pagestyle{scrheadings}  % Stil für die Seite setzen
\automark{section}  % Abschnittsnamen als Seitenbeschriftung 
\setheadsepline{.5pt}    % Kopzeile durch Linie abtrennen
\usepackage[hidelinks]{hyperref}  % Links und weitere PDF-Features
\usepackage{typearea, tikz, pgffor, verbatim}
\usetikzlibrary{calc, arrows.meta,arrows, graphs, quotes, positioning}
\usepackage[cache=false, outputdir=build]{minted}
\usemintedstyle{vs}
\usepackage{placeins} % provides \FloatBarrier

\definecolor{bg}{rgb}{0.9,0.9,0.9}
\newcommand\numberthis{\addtocounter{equation}{1}\tag{\theequation}}
\usepackage{svg}
\usepackage{siunitx}
\usepackage{physics}
\usepackage{tensor}
\renewcommand{\bra}{}
\newcommand{\kom}{}
\newcommand{\brak}{}
\renewcommand{\braket}{}
\newcommand{\intx}{}
\newcommand{\intr}{}
\newcommand{\kla}{}
\newcommand{\klam}{}
\newcommand{\klammer}{}
\newcommand{\oper}{}
\newcommand{\tecxt}{}
\newcommand{\Bra}{}
\newcommand{\Ket}{}
\renewcommand{\i}{\text{i}}
\newcommand{\e}{\text{e}}
\newcommand*{\Fourier}{\textsc{Fourier}\ }
\newcommand*{\F}{\mathcal{F}}
\newcommand*{\sinc}{\text{sinc\,}}
\newcommand{\conv}{\mathop{\scalebox{3.5}{\raisebox{-0.38ex}{\makebox[0.5\width][c]{$\ast$}}}}\limits}

\newcommand*{\titel}{A physicist's approach to the cubic equation}
\newcommand*{\autor}{Joachim Schwardt}

\hypersetup{pdfauthor={\autor}, pdftitle={\titel}}

%\titlehead{titlehead}
\subject{Number Theory}
\title{\titel}
\author{\autor}
\date{\today}

\begin{document}

%\begin{titlepage}
\maketitle  % Titel setzen
%\tableofcontents  % Inhaltsverzeichnis setzen
%\end{titlepage}

\section{Solving the quadratic}
Consider the general quadratic equation
\begin{align}
0 &= ax^2+bx+c
\end{align}
where $a\neq 0$.
Dividing through $a$ reduces the number of independent parameters and yields
\begin{align}
0 &= x^2 + px + q
\end{align} 
where $p=\frac{b}{a}$ and $q=\frac{c}{a}$.
Now there are many ways of solving this equation, but most of them get their intuition from the concept of ``completing the square''.
However, we will take a different approach.
\\
From a physicist's point of view, this is a ``hard problem'', in the sense that there is no obvious way of solving it ``directly''.
So, let us consider a special case that we \emph{can} solve.
For $q=0$, we reduce the equation to a linear one -- that is certainly easy to solve, but will not lead to further insights.
The other case of $p=0$ is more interesting, because it introduces a symmetry.
Furthermore, we can very easily solve the resulting equation
\begin{align}
0 &= x^2+q.
\end{align}
The two solutions are just $x_{1,2} = \pm \sqrt{-q}$.
Now comes the difficult step, where one has to find a way of reducing the general problem to the simple case.
The intuition here is that the $px$-term merely results in a shift of the parabola, that breaks the symmetry.
So, let us try to transform into a new set of coordinates, where the symmetry is restored.
The easiest attempt at this is a simple shift $y:=x+\Delta$.
Inserting this gives
\begin{align}
0 &= \kla\left( y-\Delta \right)^2 + p(y-\Delta) + q = \\
&= y^2 + y\kla\left( p - 2\Delta \right) + \Delta^2 - p\Delta + q.
\end{align}
We may now choose $\Delta = \frac{p}{2}$, which will eliminate the asymmetric term.
We have thus successfully reduced the general problem to one that we can solve.
We know that the solution to
\begin{align}
0 &= y^2 + q - \kla\left( \frac{p}{2} \right)^2
\end{align}
is simply $y_{1,2} = \sqrt{\kla\left( \frac{p}{2} \right)^2 - q}$.
Transforming back to our original coordinates, we have
\begin{align}
x_{1,2} &= -\frac{p}{2} \pm \sqrt{\kla\left( \frac{p}{2} \right)^2 - q}.
\end{align}
Of course, this is a somewhat lengthy derivation of a simple result -- but this technical approach will prove very useful for the more difficult cubic equation discussed in the next section.
\section{Solving the cubic}
Consider the general cubic equation
\begin{align}
0 &= ax^3+bx^2+cx+d
\end{align}
where $a\neq 0$.
Let us just be as optimistic as possible, and repeat the steps we took to solve the quadratic.
That came down to dividing through the leading coefficient and shifting the coordinate system to $y:=x+\Delta$.
Then we find
\begin{align}
0 &= \kla\left( y-\Delta \right)^3 + \frac{b}{a} \kla\left( y-\Delta \right)^2 + \frac{c}{a}\kla\left( y-\Delta \right) + \frac{d}{a} = \\
&= y^3 + y^2\kla\left( \frac{b}{a} - 3\Delta \right) + y\kla\left( \frac{c}{a} - 2\Delta\frac{b}{a} + 3\Delta^2 \right) +\frac{d}{a} - \frac{c}{a}\Delta + \frac{b}{a}\Delta^2 - \Delta^3.
\end{align}
Choosing $\Delta = \frac{b}{3a}$ will eliminate the quadratic term and leave us with
\begin{align}
0 &= y^3 + py+q
\end{align}
where we defined the coefficients $p := -\frac{b^2}{3a^2}$ and $q := \frac{d}{a} - \frac{bc}{3a^2} + \frac{8b^3}{27a^3}$.
Unfortunately, this is still a hard problem. 
But again, being very optimistic, maybe repeating the approach can help us solve the equation.
After all, we started with one more term than before, so maybe we need an additional shift $z := y + u$
\begin{align}
0 &= (z-u)^3 + p(z-u)+q = \\
&= z^3 - 3z^2u + 3zu^2 - u^3 + p(z-u) + q = \\
&= z^3 - u^3 + (p- 3zu)(z-u) + q.
\end{align}
And now we choose $u$ in such a way that removes the ``problematic'' term, i.e. $u = \frac{p}{3z}$.
Then we are left with
\begin{align}
0 &= z^3 - \frac{p^3}{27z^3} + q,\\
\Rightarrow\ \ \ 0 &= (z^3)^2 + qz^3 - \kla\left( \frac{p}{3} \right)^3.
\end{align}
This is a quadratic equation in $z^3$!
The solution is
\begin{align}
z_{1,2}^3 &= -\frac{q}{2} \pm \sqrt{\kla\left( \frac{q}{2} \right)^2 + \kla\left( \frac{p}{3} \right)^3}.
\end{align}
Since $u^3 = z^3 + q$ and $y=z-u$, we find the first solution to our cubic equation to be
\begin{align}
y_{1} &= \sqrt[3]{-\frac{q}{2} + \sqrt{\kla\left( \frac{q}{2} \right)^2 + \kla\left( \frac{p}{3} \right)^3}} - \sqrt[3]{\frac{q}{2} + \sqrt{\kla\left( \frac{q}{2} \right)^2 + \kla\left( \frac{p}{3} \right)^3}} = \\
&= \sqrt[3]{-\frac{q}{2} + \sqrt{\kla\left( \frac{q}{2} \right)^2 + \kla\left( \frac{p}{3} \right)^3}} + \sqrt[3]{-\frac{q}{2} - \sqrt{\kla\left( \frac{q}{2} \right)^2 + \kla\left( \frac{p}{3} \right)^3}}.
\end{align}
Note that choosing the negative branch in the quadratic equation leads to the same result.
%Factoring the polynomial as $y^3+py+q=(y-y_1)(y-y_2)(y-y_3)$, we find $q=-y_1y_2y_3$, $p = y_1y_2+y_1y_3+y_2y_3$ and $0 = y_1+y_2+y_3$.
%We may then write
%\begin{align}
%0 &= y_1 + y_2 - \frac{q}{y_1y_2},\\
%\Rightarrow\ \ \ 0 &= y_2^2 + y_1y_2 - \frac{q}{y_1}.
%\end{align}
%Note that
%\begin{align}
%y_1^3 &= -q -py_1
%\end{align}
%Thus the second and third solution are
%\begin{align}
%y_{2,3} &= -\frac{y_1}{2} \pm \sqrt{\frac{y_1^2}{4} + \frac{q}{y_1}} = \frac{1}{\sqrt{y_1}} \kla\left( -\frac{y_1^{3/2}}{2} \pm \sqrt{\frac{y_1^3}{4} + q} \right) = \\
%&= y_1\kla\left( \sqrt{\frac{1}{4} + \frac{q}{-q + z}} \right)•
%\end{align}
%Thus the third solution is 
%\begin{align}
%y_3 &= -y_1-y_2 =  \sqrt[3]{\frac{q}{2} + \sqrt{\kla\left( \frac{q}{2} \right)^2 + \kla\left( \frac{p}{3} \right)^3}} - \sqrt[3]{-\frac{q}{2} + \sqrt{\kla\left( \frac{q}{2} \right)^2 + \kla\left( \frac{p}{3} \right)^3}},\\
%&\sqrt[3]{\frac{q}{2} - \sqrt{\kla\left( \frac{q}{2} \right)^2 + \kla\left( \frac{p}{3} \right)^3}} - \sqrt[3]{-\frac{q}{2} - \sqrt{\kla\left( \frac{q}{2} \right)^2 + \kla\left( \frac{p}{3} \right)^3}}
%\end{align}
Figuring out the other two solutions is a little cumbersome, but the difficult step of actually finding a general solution is completed.
\section{Solving the quartic}
Consider the general quartic equation
\begin{align}
0 &= ax^4+bx^3cx^2+dx+e
\end{align}
where $a\neq 0$.
What worked twice might work again. 
We will divide through $a$ and shift to a new coordinate $y := x + \Delta$
\begin{align*}
0 &= (y-\Delta)^4 + \frac{b}{a} (y-\Delta)^3 + \frac{c}{a} (y-\Delta)^2 + \frac{d}{a} (y-\Delta) + \frac{e}{a} = \numberthis \\
&= y^4 + y^3\kla\left( \frac{b}{a} - 4\Delta \right) + y^2\kla\left( \frac{c}{a} - \frac{3b}{a}\Delta + 6\Delta^2 \right) + \\
&\hspace{0.4cm} + y\kla\left( \frac{d}{a} - \frac{2c}{a}\Delta + \frac{3b}{a}\Delta^2 - 4\Delta^3 \right) + \frac{e}{a} - \frac{d}{a}\Delta + \frac{c}{a} \Delta^2 - \frac{b}{a} \Delta^3 + \Delta^4. \numberthis 
\end{align*}
Choosing $\Delta = \frac{b}{4a}$ eliminates the cubic term.
Then we are left with
\begin{align}
0 &= y^4 + py^2 + qy + r.
\end{align}
where we defined $p := \frac{c}{a} - \frac{3b^2}{8a^2}$, $q := \frac{d}{a} - \frac{bc}{2a^2} + \frac{b^3}{8a^3}$ and $r := \frac{e}{a} - \frac{bd}{4a^2} + \frac{b^2c}{16a^3} - \frac{3b^4}{256a^4}$.
As before, let $z := y + u$.
%Then we find
%\begin{align}
%0 &= (z-u)^4 + p(z-u)^2 + q(z-u) + r = \\
%&= z^4 + u^4 - 4zu(z^2+u^2) + 6z^2u^2 + p(z^2+u^2-2zu) + q(z-u) + r = \\
%&= z^4+u^4 + (p-4zu)(z^2+u^2) - 2zu(p - 3zu) + q(z-u)+r.
%\end{align}
%Choosing $u = \frac{p}{4z}$ leads to
%\begin{align}
%0 &= z^4 + u^4 - 2z^2u^2 + q(z-u)+r = \\
%&= (z^2 - u^2)^2 + q(z-u) + r = \\
%&= \kla\left( z^2 - \frac{p^2}{16z^2} \right)^2 + q\kla\left( z - \frac{p}{4z} \right) + r,\\
%\Rightarrow\ \ \ 0 &= z^8 - \frac{p^2}{16} z^4 + \frac{p^2}{256} + qz^5 - \frac{qp}{4} z^3 + rz^4.
%\end{align}
%%%% NEW ATTEMPT
%Now we will need to eliminate two terms instead of one, so it should not be too surprising that we introduce not one but two new parameters.
%Let $z := y + u + v$.
%Then we find
%\begin{align*}
%0 &= (z-u-v)^4 + p(z-u-v)^2 + q(z-u-v) + r = \\
%&= z^4 + (u+v)^4 - 4z(u+v)\kla\left( z^2+(u+v)^2 \right) + 6z^2(u+v)^2 + \\
%&\hspace{0.4cm} + p \kla\left( z^2+(u+v)^2-2z(u+v) \right) + q(z-u-v) + r = \\
%&= z^4 + u^4 + v^4 + 4uv(u^2+v^2) + 6u^2v^2 + \\
%&\hspace{0.4cm} + (p-4zu - 4zv)\kla\left( z^2 + u^2 + v^2 +2uv \right) - \\
%&\hspace{0.4cm} - 2z(u+v) \kla\left( p - 3zu - 3zv \right) + q(z-u-v)+r .
%\end{align*}
%%%% NEW ATTEMPT
Then we find
\begin{align}
0 &= (z-u)^4 + p(z-u)^2 + q(z-u) + r = \\
&= z^4 + u^4 - 4z^3u + 6z^2u^2 - 4zu^3 + p(z^2+u^2-2zu) + q(z-u) + r.
\end{align}
Now we choose $u$ in such a way that $0 = z^4 + u^4 + r$ remains.
This demands
\begin{align}
0 &= 6z^2u^2 - 4z^3u - 4zu^3 + pz^2 + pu^2 - 2pzu + qz - qu + r = \\
&= -4zu^3 + u^2\kla\left( 6z^2 + p \right) + u\kla\left( -4z^3 - 2pz - q \right) + pz^2+qz+r,
\end{align}
which is just a cubic equation in $u$ -- we already know how to solve this.
Now we can turn to solving the quartic equation, since
\begin{align}
0 &= z^4 + u^4... \text{this is where it gets diffucult, the attempt doesnt quite work}
\end{align}



%\begin{thebibliography}{9}
%\bibitem{example} description, \url{https://www.exmapleurl}, day.\,month~2021.
%\end{thebibliography}

\end{document}